%---------------------------------------------------------------------------%
%->> Frontmatter
%-
%---------------------------------------------------------------------------------------------------------------
%-> 摘要
%-
%- 摘要内容要求在800～1200字，应简要说明本论文的目的、方法、结果和结论。
%- 要突出论文的创新之处。语言力求精炼、准确。在本页的最下方另起一行，
%- 注明本文的关键词(3-5个)。
%摘要是以提供文献内容梗概为目的，不加评论和补充解释，简明、确切地记述文献重要内容的短文。其要素一般包括：
% -(1)目的——研究、研制、调查等的前提、目的和任务，所涉及的主要范围；
% -(2)方法——所用的原理、理论、条件、对象、材料、工艺、结构、手段、装备、程序等；
% -(3)结果——实验的、研究的结果，数据，被确定的关系，观察结果，得到的效果，性能等；
% -(4)结论——结果的分析、研究、比较、评价、应用，提出的问题，今后的课题，假设，启发，建议，预测等； 
%
%写摘要时不得简单地重复题名中已有的信息，要排除在本学科领域中已成常识的内容，要用第三人称的写法。
%应采用“对……进行了研究”、“报告了……现状”、“进行了……调查”等记述方法，不使用“本文”、“作者”等作为主语。
%摘要的第一句不要与题目重复；取消或减少背景信息，只表示新情况、新内容；不说空洞的词句，
%如“本文所讨论的工作是对过去×××的一个极大地改进”、“本工作首次实现了……”、“经检索尚未发现与本文类似的工作”等；
%此外，作者的打算及未来的计划不能纳入摘要。
%
%列出3～5个关键词，以关键字与课题间关联性由强到弱顺序排列，关键词之间用逗号相隔，结束处不用标点符号。

\begin{abstract}
    
指纹在生物识别应用中的形式最为广泛，具有唯一性和不变性。指纹增强的主要任务是恢复指纹脊线的结构缺陷，其目的是通过改善指纹脊线质量来提高指纹特征提取的精度，进而提高指纹识别的精度。

在系统研究指纹增强算法的基础上，介绍了指纹增强的研究背景及现状，论述了指纹增强的概念及指纹增强的各种实现方法。除空间域增强和频域增强外，重点介绍了Gabor滤波器在指纹增强技术中的应用，并且在此基础上提出了改进算法。仿真显示Gabor滤波法在指纹增强技术中具有很好的带通特性，也具有很好的方向和频率选择特性。因此，用Gabor滤波器对原始指纹图像进行滤波可以有效地去除脊线方向的噪声，保存真实的脊线结构。Log-Gabor滤波法作为改进算法，弥补了Gabor滤波法的缺陷，改善指纹图像的滤波效果。

\keywords{指纹增强，脊线方向，脊线频率，Gabor滤波}
\end{abstract}

%英文摘要上方应有题目，内容与中文摘要相同。
%在英文题目下面第一行写研究生姓名，专业名称用括弧括起置于姓名之后，
%研究生姓名下面一行写导师姓名，格式为Directed by...。
%最下方一行为英文关键词(Keywords 3-5个)
%英文摘要撰写注意事项：
%(1) 尽量取消不必要的字句：如“It is reported…”，“The author discusses…”，“In this paper，”。 
%(3) 采用短句叙述，但要避免句形单调；目的、方法、结果用过去时态，结论用一般现在时态；可用动词的情况应尽量避免用
%动词的名词形式；避免使用一长串形容词或名词来修饰名词；注意冠词用法，不要误用、滥用或随便省略冠词。
%(4) 语言要精炼，多用简短、词义清楚、熟悉的词；避免使用文学性的描述手法撰写文摘。
%(5) 对已经为大众所熟悉的缩写词可直接使用，对于那些仅为同行所熟悉的缩略语，应在题目、摘要或关键词中至少出现一次全称。
\begin{abstract*}
In biometric identification, the most widely used form was the fingerprint, which has unique and invariable property. The main task of fingerprint enhancement was to restore the structural defects of its ridge. The ultimate goal was to improve the accuracy of fingerprint feature extraction by improving the quality of the ridge, leading to improving the accuracy of fingerprint identification.

Based on the fingerprint enhancement algorithms, the paper introduced the research background, actuality of fingerprint enhancement, the concept of fingerprint enhancement and various implementation methods. Other than spatial and frequency domain enhancement, the paper focused on introducing the Gabor filter in fingerprint image enhancement techniques and putting forward some improved algorithms. The Gabor filtering has nice band pass ability in fingerprint enhancement from simulation results, with nice direction and frequency selectivity. Thus, Gabor filtering can effectively remove the noise of the ridge along its direction, also save the true ridge structure. As an improved algorithm, Log-Gabor filter which can make up the defect of Gabor filter improves the final effect of filtering.

\keywords*{Fingerprint enhancement, Ridge orientation, Ridge frequency, Gabor filtering}
\end{abstract*}
